\scanpage{25}
\diagram{
\begin{tikzcd}[column sep=small,row sep=small]
\dot T\cap E\simeq\dot T_0\times I
\arrow[dr,phantom,"\subset" sloped] & \\
\deleted{\unsure{et où}} & T\cap E\simeq T_0\times I \\
Y\cap E\unsure{\}}\simeq Y_0\cap E
\arrow[ur,phantom,"\subset" sloped] &
\end{tikzcd}}
où
\[
\dot T_0=\dot T\cap X_0,\qquad
T_0=T\cap X_0,\qquad
Y_0=Y\cap X_0.
\]

Et si, symétriquement,
\deleted{$\unsure{T\cap E},T\cap\dot E,X_0\quad\unsure{T\cap E}$}

b) $E\cap T\ (=T\cap E)$, $\dot E\cap T\deleted{\simeq T_0\times\{1\}}$,
$X_0\cap T\ (=T\cap X_0=T_0)$ sont des
sous-espaces côniques de $T$, donc en particulier
stables par $T$-homothéties. \deleted{Le fait que} Noter que
les trois espaces précédents s'identifient, grâce à a),
à $T_0\times I$, $T_0\times\{1\}$, $T_0\times\{0\}$. On suppose enfin

c) Les homothéties \deleted{par} $\lambda_E$ et $\lambda_T$ commutent
dans $T\cap E$, ce qui implique que l'application
\deleted{$\unsure{\lambda}$} induite par $\lambda_T$ ($\lambda\in[0,1]$) sur
$T_0\times\{0\}$ et $T_0\times\{1\}$
s'identifie à une $=$ \deleted{homothétie} opération de
$T_0-T_0$, \deleted{et que} soit $\lambda_{T_0}$, et que l'opération
$\lambda_T$ sur $T\cap E\simeq T_0\times I$ s'identifie à
\deleted{$\unsure{\lambda_{T_0}}$} $\lambda_{T_0}\times\operatorname{id}_I$.

Donc les conditions a) b) c) équivalent à celle-ci :

\uline{Condition de compatibilité} $\exists$ homéomorphisme
\[
T\cap E\simeq T_0\times I
\qquad\text{(compatible avec }E\simeq X_0\times I
\]
\deleted{compatible aux structures}
induisant
\[
\dot T\cap E\simeq\dot T_0\times I,\qquad Y\cap E\simeq Y_0\times I
\]
tels que la structure cônique de $T$ sur $Y$
\unsure{s'identifie} sur $\overline T_E$ sur $Y_E$
\deleted{avec \unsure{de} la} \inserted{une}
\uline{structure cônique produit} \inserted{par $I$}
d'une structure cônique de
$T_0$ sur $Y_0\subset\overline T_0$, \deleted{par $I$}.
\margin{24}
